\chapter*{Résumé} 
\addcontentsline{toc}{chapter}{Résumé}

Ce mémoire présente le développement de MindCare-AI, une solution innovante dédiée au suivi et à l’amélioration du bien-être mental. Réalisé au sein de la société Mobelite, ce projet vise à répondre aux défis majeurs de la santé mentale en proposant une application mobile intelligente, accessible et personnalisée. MindCare-AI combine l’analyse de données physiologiques (par photopléthysmographie à distance) et l’analyse de textes pour évaluer l’état émotionnel et physiologique des utilisateurs.
Le système intègre plusieurs modules complémentaires : acquisition et traitement de signaux vitaux, reconnaissance automatique des émotions à partir du langage naturel, extraction de caractéristiques pertinentes, et classification à l’aide de modèles d’intelligence artificielle. Un chatbot thérapeutique interactif vient enrichir l’expérience utilisateur en offrant un accompagnement et des conseils personnalisés en temps réel.
La méthodologie adoptée s’appuie sur une démarche agile pour le développement logiciel et sur des techniques avancées de data science pour la préparation, le prétraitement et l’analyse des données, le tout automatisé dans des notebooks Jupyter pour garantir la reproductibilité. Les résultats expérimentaux obtenus démontrent la pertinence et l’efficacité de l’approche proposée pour le suivi continu, non invasif et personnalisé de la santé mentale, ouvrant la voie à de nouvelles perspectives en e-santé et en prévention psychologique.